\subsection{From construction to recognition}

The geometric definitions above answer the question \emph{what the curves are}.
A different problem begins when an equation is given first:

\begin{center}
\emph{Which geometric curve, if any, is hidden in this algebraic condition?}
\end{center}

When no mixed term is present, completing squares already reveals the basic
geometric alternatives. Squared terms with the same sign suggest a bounded
ellipse-type set, opposite signs suggest hyperbolic behaviour, and the absence
of one quadratic direction suggests a parabola. Degenerate cases may instead
produce points, lines, or no real points.

The mixed term $Bxy$ indicates that the chosen coordinate axes may not follow the
natural symmetry directions of the curve. A rotation of coordinates can remove
that term. The next section develops the general second-degree equation and the
rotation-invariant discriminant used to classify it.

\begin{remarkbox}
The purpose here is only to motivate the recognition problem. The precise
classification criterion and its exceptional cases are stated once, in the
algebraic classification that follows.
\end{remarkbox}
